% RESEARCH_DRAFT_SHA256: 73a5eb464dbbdb563a85555f881f65387e0a10217d767c609cf3bc3dec618544
% Layout and source-reference wrappers expanded for the web renderer.
\documentclass[11pt]{article}
\usepackage[margin=25mm]{geometry}
\usepackage{amsthm}
\usepackage{fontspec,unicode-math}
\setmainfont{Latin Modern Roman}
\setmathfont{Latin Modern Math}
\usepackage{amsmath,amssymb,mathtools,xurl,hyperref}
\hypersetup{colorlinks=true,urlcolor=blue,linkcolor=blue}
\setcounter{secnumdepth}{0}
\setlength{\parindent}{0pt}
\setlength{\parskip}{5pt}
\setlength{\emergencystretch}{3em}
\newcommand{\sref}[2]{\hyperlink{source:#1:#2}{[S#2]}}
\newcommand{\cataloguescope}[1]{\paragraph{Source record.}#1}
\newcommand{\floor}[1]{\left\lfloor#1\right\rfloor}
\newtheorem{proposition}{Proposition}
\newtheorem{lemma}[proposition]{Lemma}
\begin{document}
\paragraph{Research manuscript.} OpenAI \texttt{gpt-6-astra}, reasoning \texttt{ultra}; 11 October 2026.
\paragraph{Preserved source body.} The complete catalogue mathematical text and sources follow with layout and source links expanded, including source qualifications and its historical preparation notes. New research begins at the section headed Research attempt.

% UnsolvedMath ID 30005631; source catalogue code OWR-14297738-004
\section{Convergence to Implicit Geometry under Class Imbalance}\label{30005631}
{\small UnsolvedMath ID 30005631 · Computer Science}
\subsection{Definitions and mathematical statement}

\paragraph{Shared notation.}Standard number systems are $\mathbb N,\mathbb Z,\mathbb Q,\mathbb R,\mathbb C$. Specialized notation and index ranges follow the cited source and the definitions below.


Class imbalance means that class frequencies differ. An optimization trajectory can select a limiting parameter or representation geometry even among solutions with equal training loss; a convergence rate requires a specified distance and time or iteration scale.
\paragraph{Record scope and qualifications.}
The catalogue title is retained as metadata. Named certificates, contracts, bounds or frameworks in that title are not asserted as results of the original source extract.
Architecture, data, loss, initialization, and the geometry being approached are not specified.

Source-record limitation: The preferred implicit geometry, algorithm, data model and error metric for convergence rates are unspecified; the question remains a broad research direction.
\paragraph{Statement recorded in the supplied source.}
Quantify how class imbalance affects the rates at which first-order training methods converge to their asymptotically preferred implicit geometry, and determine how architecture, data, and loss design affect that convergence. \href{https://doi.org/10.4171/owr/2023/41}{[S2]}
\paragraph{Source status.}
The supplied catalogue labels this source record partially solved; that metadata was checked against the supplied snapshot on 2026-10-11. The label is preserved for provenance and is not an independent verification of present open status or mathematical correctness. The catalogue literature assessment is dated 2026-08-22.
\subsection{Short English statement}
Quantify how class imbalance alters the geometry and convergence of trained models.
\subsection{Sources}
\begin{itemize}
\item[S1] Ulam AI and the UnsolvedMath Contributors, UnsolvedMath: A Curated Collection of Open Mathematics Problems, record 30005631 (OWR-14297738-004). Catalogue curation and metadata: CC BY 4.0. The source extract is retained; definitions, English summaries and qualifications are editorial additions.
\url{https://huggingface.co/datasets/ulamai/UnsolvedMath}.
\item[S2] Mini-Workshop: Interpolation and Over-parameterization in Statistics and Machine Learning​ (2024), Oberwolfach Reports, DOI 10.4171/owr/2023/41, pp. 2359--2376
\url{https://doi.org/10.4171/owr/2023/41}.
\item[S3] Kini--Vakilian--Behnia--Thrampoulidis, imbalance geometry
\url{https://openreview.net/forum?id=utahaTbcHdP}.
\item Source-specific Oberwolfach Reports licences, including CC BY-SA 4.0 where applicable, are retained; see THIRD\_PARTY\_NOTICES.md.
\end{itemize}
\paragraph{Preparation and review.}
The original source and UnsolvedMath contributors retain attribution for the source extract. Codex (GPT-6; exact serving revision not exposed) assisted with definitions, English summaries, scope notes and formatting. No new human mathematical authorship or independently verified solution is asserted. Permanent source ID 30005631 and code OWR-14297738-004 checked against the complete supplied registry and current repository. Distinct upstream records, including repeated questions and broad or damaged extracts, are retained for complete coverage of the source CS category. A repeated record is not a new independent conjecture. Complete mathematical-equivalence and current-literature review has not been performed. Preparation and final review follow contribution prompts C and D, with source gaps disclosed. Structural and rendering checks do not establish a new theorem or open status.
\label{end:30005631}

\section{Research attempt}
\subsection{Literature, source context and scope}
The full original catalogue body and source qualifications are preserved above. The relevant primary workshop passages were read in the cached publisher report. Historical claims and subsequent attributed results are distinguished from the mathematical work below. A restricted example, standard lemma or known theorem is not counted as new progress on the entire original target. Missing assumptions and repeated or rhetorical extracts are not silently turned into different solved problems.

\subsection{An orthogonal logistic model separates direction and rate}
The source develops feature-learning geometries and asks about convergence under imbalance. We use a much smaller fixed-feature model whose gradient flow can be integrated exactly. Let $0<p<1$ be the fraction of one class. After multiplying each feature by its binary label, suppose the two types are $e_1$ and $e_2$. The logistic loss is
\[
L(w)=p\log(1+e^{-w_1})+(1-p)\log(1+e^{-w_2}),
\]
and gradient flow starts at $w(0)=0$. Both types are separable; their different coefficients represent the class frequencies.
\begin{proposition}
Writing $p_1=p$, $p_2=1-p$, the exact trajectory satisfies
\[
e^{w_j(t)}+w_j(t)=1+p_jt.
\]
For fixed $p\in(0,1)$, $w(t)/\|w(t)\|\to(1,1)/\sqrt2$. If $p\ne1/2$, the angle $\theta(t)$ to that direction satisfies
\[
\tan\theta(t)\sim\frac{|\log(p/(1-p))|}{2\log t}.
\]
\end{proposition}
\begin{proof}
Gradient flow gives $w'_j=p_j/(1+e^{w_j})$, so differentiating $e^{w_j}+w_j$ gives the constant $p_j$. Its initial value is one, proving the formula and uniqueness because $e^w+w$ is strictly increasing. Since $w_j\ge0$, the equation implies $w_j\le\log(1+p_jt)$. Therefore $w_j/t\to0$ and $e^{w_j}/(p_jt)=1+(1-w_j)/(p_jt)\to1$. Taking logarithms yields $w_j=\log t+\log p_j+o(1)$. The direction conclusion follows. The parallel and perpendicular coordinates relative to $(1,1)/\sqrt2$ are $(w_1+w_2)/\sqrt2$ and $(w_1-w_2)/\sqrt2$, so their absolute ratio is the tangent of the angle. Substitution proves the asymptotic when its numerator limit is nonzero. For $p=1/2$, the coordinates are identical for all $t$ and the direction is already balanced for $t>0$.
\end{proof}
The convergence is not uniform as the minority frequency goes to zero. Since $w'_j\le p_j$, reaching any minority margin $R>0$ needs at least $R/p_j$ time. Reweighting the two loss terms to have coefficients $1/2$ instead gives identical coordinate trajectories from this initialization, but that is a changed objective and does not prove improved population performance.
\subsection{Verification and remaining gaps}
The implicit equations were solved numerically by monotone bisection for several imbalances and times, verifying their residuals, angular relation and minority-margin bound. The limit takes time to infinity at fixed positive $p$; interchanging this with $p\downarrow0$ is unjustified. The model has fixed orthogonal features, not evolving embeddings or a SELI neural-collapse geometry. Thus a common limiting direction in this example does not establish architecture-independent geometry or rate. The listed OpenReview page returned a browser challenge; the primary workshop account was read, and no inaccessible paper proof is claimed checked.

\section{References and source attribution}
Primary workshop report: \url{https://doi.org/10.4171/owr/2023/41}. Relevant passages read: 2023/41 pp. 2370--2371, SELI geometry, imbalance and convergence-rate questions.

\url{https://openreview.net/forum?id=utahaTbcHdP}. Attempted primary paper landing page; browser challenge prevented reading its theorem. Cached workshop source supplies the scoped SELI discussion.

All retrievals were checked on 11 October 2026. Original question attribution remains in the preserved Sources section. External theorems are attributed and not claimed as new discoveries or fully verified merely by citation.


\paragraph{Integration review.} Codex (GPT-6; exact serving revision not exposed) reviewed the written arguments and expanded layout and source-reference wrappers. Repository metadata and the production web renderer were checked using Python 3.13 and JavaScript with MathJax 3.2.2. This model-assisted review is not independent human proof verification; no PDF compilation or formal proof checking is asserted.

\section{Completion Estimate}
\noindent\textbf{COMPLETION: 0\%}
No substantiated new progress toward the entire original target is claimed. The worked results and failed route are delimited above; a known result, illustrative model, or auxiliary obstruction is not counted as original-target completion.

\section{Final status and attribution}
\textbf{Attempt status: unresolved.} This describes this research attempt, not a new ruling on the source catalogue's status or on external literature claims.
The mathematics and this self-audit were generated by OpenAI \texttt{gpt-6-astra}, reasoning \texttt{ultra}, on 11 October 2026. Codex (GPT-6; exact serving revision not exposed) orchestrated the batch, read and reviewed its written mathematical arguments, and contributed formatting/repository review. No human mathematical authorship, independent human proof checking, or formal verification is claimed. The historical preparation paragraphs in the reproduced catalogue body concern that body only.

Contribution prompt A was applied to the complete definition and supplied source record, followed by the adversarial checks of prompt B. Local reproduction files are \texttt{batch025/checks.py}, \texttt{computations.json}, and the per-ID audit JSON. They were not public links at preparation time. Numerical checks are finite and do not prove unrestricted results. No LaTeX engine was found locally; compilation was not performed by this research agent.

\end{document}
